\documentclass[12pt]{article}

\title{An urban model of elasticity of price with respect to unit size}
\author{Tom Davidoff and Tsur Somerville}


Households have utility $U\left(c,s,\theta\right)$, where
$c$ is numeraire consumption, $\theta$ is a single index of neighbourhood
quality (e.g. distance from downtown), and $s$ is the square feet of housing
consumed at a lot at their chosen location $\theta$. $U_{c}>0$, $U_{s}>0$,
$U_{\theta}>0$, and $U_{cc}<0$.

Households vary only by income $y$. Numeraire consumption is given by:

\begin{equation}
	c = y - sp(s,\theta)+m\theta.
\end{equation}

The term $m\theta$ in consumption recognizes the possibility that there is a
direct monetary benefit to living in a better location (e.g. less commuting or
more networking opportunities). $\theta$ also enters utility directly, e.g.
through amenity or the time value of commuting shorter distances. Price per
square foot of housing $p$ is a function of both the location quality $\theta$,
and the size of the unit $s$. 


Developers can build up to density $\bar{s}$ at each location, and government
rations a small number of permits to build at $k\bar{s}$, with $k<1$. 

For households of a given income $y$, the derivative of the price per square
foot with respect to location quality $\theta$ is given by noting that utility
must not change with distance along a bid rent curve.

\section{Sorting on income}
\begin{align}
	\frac{dU}{d\theta}|_{s=\bar{s}} &= -U_{c}\left[-\bar{s}\frac{dp}{d\theta}+m\right]+U_{\theta} = 0 \\
	\frac{dp}{d\theta}|_{s=\bar{s}} &= \frac{mU_{c}+U_{\theta}}{\bar{s} U_c}\\
\frac{d^{2}p}{d\theta dy} & = \frac{mU_{cc} + U_{c\theta} -\bar{s}\frac{U_{cc}{U_c}}}{\bar{s}U_{c}}
\end{align}

Because $U_{cc}<0$, the second derivative of interest is positive (so that higher income households have steeper bid rent curves and outbid lower income households for better locations) if $m$ is suffciently small and $U_{c\theta}\geq 0$. That is, if amenity and time features of location are more important than monetary commuting costs, there will be positive sorting into location quality by income.

\section{Effect of size on price per square foot}

Suppose $k\approx 1$, so that the government zoning reform is very small. In
that case, introduction of a smaller unit size will not affect spatial sorting
by income even in the predictable case that there is sorting by income into
house size (assuming the sorting by income into location quality is strict).

The slope of price per square foot around $\bar{s}$ when there is free choice
of unit size is:

\begin{align}
	\frac{dU}{ds}|_{s=\bar{s}} & = -U_{c}\left[p(\bar{s},\theta)+\bar{s}\frac{dp}{ds}\right] + U_s = 0 \\
	\frac{dp}{ds}|_{s=\bar{s}} & = \frac{U_{s}-U_{c}p(\bar{s},\theta)}{\bar{s} U_c}.\\
	\frac{d^2 p}{d s d y} & = \frac{U_{cs} - U_{cc}\left[p(\bar{s},\theta) +\frac{1}{U_{c}}\right]}{\bar{s} U_c}.
\end{align}

As long as $U_{cs}$ is not too negative, the second derivative of price per
square foot with respect to size and income will be positive by concavity of
utility over consumption. 

Combining results, we know that for fixed housing consumption, under reasonable
conditions, there is positive sorting by income into locaiton quality, and that
the slope of willingness to pay per square foot falls with income. We thus
expect steeper declines in price per square foot with size in lower income
(worse) locations.

When the change in allowe square footage is large, as with multiplex zoning,
there could be a reshuffling of population, as small-unit low-income bid rent
curves could be steeper than large-unit high-income bid rent curves, while
large untis might dominate small units for high-income households. For this
reason, it is an empirical question whether price per square foot is, in fact,
falling more steeply in low-income locations.

\end{document}


Due to zoning, landowner developers are constrained to build $\bar{s}$ square
feet on each lot. The government auctions a small number of permits to build
homes at size $k\bar{s}$, with $k<1$. The owners of land at the locations where
the profit from building at $k\bar{s}$ minus the profit from building at
$\bar{s}$ (or not building at all, which we rule out \dots) are greatest are
allocated the permits.

For small changes in size, $k\approx 1$, we can approximate the difference in
price per square foot $p(\theta,\bar{s})$ and $p(\theta,k\bar{s})$ by the slope
of the price function with respect to size.

For homes of size $s$, indifference among locations requires


$\left[\bar{s}+\epsilon, \bar{s}-\epsilon\right]$ at each location. Housing
units at location quality $\theta$ must be built to size $\bar{s}$, but there
will be some accidental variation (e.g. through depreciation or improvements
and additions, which we do not explicitly model. Realistically, due to
grandfathering and heterogeneous zoning within neighbourhoods, the housing
stock may be fairly diverse even in the presence of currently restrictive
zoning in most locations). 

\begin{align}
	\label{eq:sFOC}
	\frac{dU}{ds} &= -U_{c}\left[p(s,\theta)+s\frac{dp}{ds}\right] + U_s \\
	\frac{dp}{ds} = \frac{U_s}{s U_c} - \frac{p(s,\theta)}{s}
\end{align}

There is no housing supply choice over $s$ to enforce equalization of price per square foot within locations.

For any household indifferent in some continuous range of $\theta$, the slope of $p$ in $\theta$ given by the ratio of the marginal utility of numeraire consumption to the net marginal benefit of increasing location quality:

\begin{equation}
	\label{eq:bidRent}
	\frac{d p}{d \theta} = \frac{mu'+U_{\theta}}{s u'}.
\end{equation}

The derivative of the ``bid-rent'' \eqref{eq:bidRent} slope with respect to income will be given by:

\begin{equation}
	\frac{d^{2}p}{d\theta dy} = \frac{mU_{cc} + U_{c\theta} + \frac{ds}{dy}p\left[mU_{cc} + U_{c\theta}\right] - \frac{sU_{cc}\left[1+\frac{ds}{dy}p\right]c{ds}{dy}U_{c}}{sU_{c}}}{sU_{c}}.
	%\frac{d^{2} p}{d \theta dy} = \frac{mU_{cc} + U_{c\theta} + \frac{ds}{dy}\left[mU_{cc} + U_{c\theta}\right] - \frac{sU_{cc}+\frac{ds}{dy}U_{c}}{sU_{c}}}{sU_{c}}.
\end{equation}
\end{document}

With negligible variation in housing consumption due to the $s$ mandate, the $\frac{ds}{dy}$ terms will be small. Then, if monetary benefits of better location are small and other amenity features are not less valued at high levels of consumption, the dominant effect will be the relative importance of money for low income households $\frac{-U_{cc}}{U_{c}}>0$.

In the small range of housing variation, the price per square foot of housing $s$ will vary within locations will come from indifference between moving to a slightly larger unit:

\begin{equation}
	\frac{dU}{ds} &= -U_{c}\left[p(s,\theta)+\frac{dp}{ds}\right] + U_s = 0\\
\end{equation}

Totally differentiating with respect to $\theta$ gives:

\begin{equation}
	0 = -U_{c\theta}\left[p(s,\theta)+\frac{dp}{ds}\right]-U_{cc}\left[
\end{equation}

\end{document}


The bid-rent function at $\theta$ has slope given by the ration $\frac\frac{{dV}{d\theta}}{\frac{dV}{dp}}$:

\begin{equation}
	\frac{d p}{d \theta} = \frac{1}{s w'(y-sp(s,\theta))}
\end{equation}

The derivative of $w'$ with respect to $y$ is negative, and hence:

\begin{equation}
	\frac{d^2 p}{d \theta d y} = -\frac{w''(y-sp(s,\theta))}{s w'(y-sp(s,\theta))^2} >0
\end{equation}

Thus higher-wealth households have steeper bid-rent curves and outbid lower
wealth households and locate in the higher-$\theta$ neighborhoods.

Within a neighborhood, the price

\end{document}


\usepackage[margin=1in]{geometry}
\usepackage{amsmath}
\usepackage{setspace}

\begin{document}

\documentclass[12pt]{article}
\usepackage{amsmath,amssymb}
\usepackage[margin=1in]{geometry}
\usepackage{setspace}
\onehalfspacing

\title{An Urban Model of the Elasticity of Price per Square Foot with Respect to Unit Size}
\author{Tom Davidoff and Tsur Somerville}

\begin{document}
\maketitle

\section{Setup}

Households differ only by income $y$ and have utility $U(c,s,\theta)$, where
$c$ is numeraire consumption, $s$ is housing consumption (square feet), and
$\theta$ indexes neighbourhood quality. Natural assumptions are $U_c>0$,
$U_s>0$, $U_\theta>0$, and $U_{cc}<0$. $\theta$ might capture distance from
downtown or other locational amenities. Locations need to be rankable on this
single index but need not be ordered geographically.

Housing is priced at $p(s,\theta)$ per square foot, allowing for nonlinear pricing in unit size. The budget constraint is:
\begin{equation}
c = y - s,p(s,\theta) + m\theta,
\end{equation}
where $m\theta$ captures monetary benefits of location (e.g.\ commuting savings). Location also enters utility directly, e.g. through amenity or time saved commuting.

Developers are  build up to a zoning limit $\bar{s}$ at each location. We analyze local behavior around $\bar{s}$, where zoning binds.

\section{Sorting by Income into Location Quality}

Along a bid-rent curve for households consuming $s=\bar{s}$, utility is constant in $\theta$:
\begin{equation}
\frac{dU}{d\theta}\bigg|{s=\bar{s}} 
= U_c\left(m - \bar{s}\frac{dp}{d\theta}\right) + U\theta = 0.
\end{equation}
Solving,
\begin{equation}
\frac{dp}{d\theta}\bigg|{s=\bar{s}} 
= \frac{U\theta + mU_c}{\bar{s}U_c} > 0.
\end{equation}

To determine sorting, differentiate with respect to income $y$, noting that $y$ enters through $c$. Applying the quotient rule:
\begin{equation}
\frac{d^2 p}{d\theta, dy}\bigg|{s=\bar{s}} 
= \frac{U{c\theta} U_c - U_\theta U_{cc}}{\bar{s} U_c^2}.
\end{equation}

Since $U_c>0$, $U_\theta>0$, and $U_{cc}<0$, the second term in the numerator is positive. Thus, the cross-partial is positive provided $U_{c\theta}$ is not too negative. In particular, if $U_{c\theta} \ge 0$, higher-income households have steeper bid-rent curves and sort into higher-quality locations. Notably, the commuting-cost parameter $m$ does not affect this sorting result.

\section{Unit Size and Price per Square Foot}

When households choose unit size freely, the first-order condition is:
\begin{equation}
U_s = U_c\left[p(s,\theta) + s\frac{dp}{ds}\right].
\end{equation}
This equates the marginal utility of space to the marginal expenditure required to acquire it. If price per square foot were constant in $s$, this reduces to the familiar condition $U_s = U_c p$. Nonlinear pricing ($dp/ds \neq 0$) reflects supply-side factors such as fixed costs or zoning constraints.

Evaluating around $s=\bar{s}$, the slope of price per square foot with respect to unit size is:
\begin{equation}
\frac{dp}{ds}\bigg|{s=\bar{s}} 
= \frac{U_s - U_c p}{\bar{s}U_c}.
\end{equation}

\subsection{Income and the Size Gradient}

Differentiating with respect to income yields:
\begin{equation}
\frac{d^2 p}{ds, dy}\bigg|{s=\bar{s}} 
= \frac{(U_{cs} - U_{cc}p)U_c - (U_s - U_c p)U_{cc}}{\bar{s} U_c^2}.
\end{equation}

At an interior optimum, $U_s = U_c p$, so this simplifies to:
\begin{equation}
\frac{d^2 p}{ds, dy}\bigg|{s=\bar{s}} 
= \frac{U{cs} - U_{cc}p}{\bar{s} U_c}.
\end{equation}

Because $U_{cc}<0$ and $p>0$, the term $-U_{cc}p$ is positive. Thus, the cross-partial is positive provided $U_{cs}$ is not too negative. Intuitively, higher-income households have lower marginal utility of consumption and are therefore less sensitive to marginal price changes, flattening the relationship between price per square foot and unit size.

\section{Implications}

Two results follow:

\begin{enumerate}
\item Higher-income households sort into higher-quality locations.
\item The decline of price per square foot with unit size is less steep for higher-income households.
\end{enumerate}

Since higher-income households occupy better locations, these results imply:

\begin{quote}
\textbf{Price per square foot declines more steeply with unit size in lower-income neighbourhoods.}
\end{quote}

This result is local, applying near the zoning-constrained size $\bar{s}$. Large changes in allowable density may alter sorting patterns and therefore the spatial gradient of price per square foot, making the empirical relationship ambiguous in such cases.

\end{document}


